\section{结论与展望}
\label{sec:conclusion}

我们展示了通过 LaMET 并结合混合比值重正化计算得到的首个核子胶子 PDF，采用两个重于物理值的π介子质量（$690$ 与 $310$ MeV），以及两种规范链涂抹技术：流时间 $3a^2$ 的 Wilson 流与 5 步超立方涂抹。
基于迈向 LaMET 胶子 PDF 的先前工作~\cite{Good:2024iur}，我们聚焦于单一胶子算符 $O^{(3)}$（定义于文献~\cite{Good:2024iur} 式 6），该算符已被用于采用赝 PDF 方法的非极化胶子研究~\cite{Fan:2020cpa,Fan:2021bcr,HadStruc:2021wmh,Salas-Chavira:2021wui,Fan:2022kcb,Delmar:2023agv,Good:2023ecp}。
我们计算了该算符混合重正化所需的单圈 Wilson 系数以及相应的单圈匹配核。
利用这些 Wilson 系数，我们发现 $\delta m+m_0$ 在π介子与核子之间保持一致，无论轻夸克质量还是奇异夸克质量——这与文献~\cite{Good:2024iur} 中 $O^{(1)}$ 与 $O^{(2)}$ 算符的先前结果形成对照。
我们对每种π介子质量与涂抹方式，在强子推动动量 $P_z = 1.71$~GeV 处计算了核子胶子 PDF。
我们发现 Wilson3 涂抹显著影响物理，给出的 PDF 在 $x \gtrsim 0.6$ 附近为负，而噪声更大的 HYP5 PDF 则不会明显低于零。
我们将两种π介子质量下的 HYP5 核子胶子 PDF 与 CT18~\cite{Hou:2019efy}、JAM24~\cite{Anderson:2024evk}、CJ22~\cite{Cerutti:2025yji} 全局拟合作了小样本比较。
我们发现与 CT18 PDF 在整个区域符合良好，而在 $x\approx 0.5$ 附近与 CJ22 及 JAM24 的一致性较差，这表明我们的 PDF 完全落在全局拟合结果的范围之内。
\UPDATE{格点 PDF 最终的误差棒并未估计为完成傅里叶变换所作外推的系统误差，而展示这类外推正是这些首批原理性验证结果的核心挑战所在。
贝叶斯与机器学习技术~\cite{Chowdhury:2024ymm,Dutrieux:2024rem,Dutrieux:2025jed} 可被移植过来，以扩展简单的外推模型，给出更严格的误差量化。}
这项研究有力地证明了用大动量有效理论计算胶子 PDF 的可行性——长期以来人们认为这难以企及。
我们相信，本工作为投入时间与资源、以降噪技术取代大π介子质量与涂抹（从而减少物理效应）提供了强有力的动机，以便从 LaMET 得到更可靠的胶子 PDF。
我们计划改进当前的计算：在多个格距上使用带零动量推动矩阵元的自重正化方案~\cite{LatticePartonLPC:2021gpi}。
\UPDATE{这将使我们能更好地理解离散化与涂抹效应，并通过取物理连续极限改善胶子 PDF 的系统误差。}
随着误差变小，还需要理解标度变化对 LaMET 胶子 PDF 的影响。
